\section{Windows from finitely many primes}\label{wt:sec:work}

\begin{resultpar}\label{wt:res:law}
Let $\mathcal P_S$ be the prime part restricted to the prime powers of a finite set of primes $S$ (all powers of its primes included), and let $t_n(S)$ be the first point where $(-1)^n(\mathcal A-\mathcal P_S)^{(n)}$ fails to be positive. Adding primes in increasing order adds the atoms $\log2,\log4,\log8,\log3,\dots$ in groups, and $\mathcal A-\mathcal P_S=W+\mathcal P_{\mathrm{rest}}$ differs from $W$ by the omitted prime powers. The first omitted one, $q$, contributes about $\Lambda(q)q^{-1/2}g_t(\log q)$, of exponential size $e^{-(\log q)^2/4t}$, while $W$ and its derivatives decay like $e^{-\gamma_1^2t}$. The two meet at
\begin{equation}\label{wt:eq:law}
t(S)\approx\frac{\log q}{2\gamma_1},\qquad q=\text{the smallest prime power not generated by }S.
\end{equation}
This is a heuristic balance of exponents: coefficients, derivative polynomials, the other zeros and the other omitted prime powers contribute corrections with no uniform bound, so \eqref{wt:eq:law} is an empirical law for this truncation scheme and not a theorem. Read as an uncertainty principle, a Gaussian of width $\sqrt{2t}$ in $u=\log n$ resolves the first missing atom when its dual width resolves the lowest zero.\end{resultpar}

Table~\ref{wt:tab:pbp} tests \eqref{wt:eq:law}, with $t_n(S)$ found by an upward scan to the first failure followed by bisection (computed, 35 and 45 digits, scanned to $0.8$).

\begin{table}[h]\centering\small
\fitwidth{\begin{tabular}{lccccc}
\toprule
primes $S$ & first omitted $q$ & $\log q/2\gamma_1$ & $t_1(S)$ & $t_2(S)$ & $t_3(S)$\\
\midrule
none & 2 & 0.0245 & 0.0253 & 0.0510 & 0.0761\\
2 & 3 & 0.0389 & 0.0394 & 0.146 & 0.223\\
2, 3 & 5 & 0.0569 & 0.0575 & 0.476 & 0.0608\\
$\le 5$ & 7 & 0.0688 & 0.0696 & $>0.8$ & 0.0720\\
$\le 7$ & 11 & 0.0848 & 0.0856 & $>0.8$ & 0.0872\\
$\le 13$ & 17 & 0.1002 & 0.1010 & $>0.8$ & 0.1022\\
$\le 31$ & 37 & 0.1277 & 0.1289 & $>0.8$ & 0.1297\\
$\le 97$ & 101 & 0.1633 & 0.1635 & $>0.8$ & 0.1641\\
\bottomrule
\end{tabular}}
\caption{End of the positivity window of order $n$ when only the primes in $S$ are kept (computed). The entries $>0.8$ mean no failure up to the end of the scan.}\label{wt:tab:pbp}
\end{table}

Three facts emerge.
\begin{enumerate}
\item The law \eqref{wt:eq:law} holds for the order one window. Its error, as a percentage of $t_1(S)$, is $3.2\%$ for $S$ empty ($q=2$), between $0.8\%$ and $1.3\%$ for the next six rows (about 1\% from $q=5$ on), and $0.17\%$ at $q=101$.
\item Among the orders $1,2,3$ tested, order $1$ always fails first. A mechanism for the sign: each omitted term $g_t(\log q)$ is increasing in $t$ for $t<(\log q)^2/2$, since $\partial_t\log g_t(u)=u^2/4t^2-1/2t$, so that $-W'\ge-(\mathcal A-\mathcal P_S)'$ there (Proposition~\ref{wt:prop:cert}(1)). For the higher derivatives $\partial_t^ng_t(u)=(-1)^nn!\,t^{-n}g_t(u)L_n^{(-1/2)}(u^2/4t)$ is positive once $u^2/4t$ exceeds the largest zero of the Laguerre polynomial $L_n^{(-1/2)}$, a threshold that grows with $n$; for a finite set of orders the omitted primes then subtract from the odd order conditions and add to the even ones on a small common interval. This accounts for the parity effect in the table ($t_2>0.8$ once the primes up to 5 are in) but not for a ranking of order 1 against every higher order, and for $S$ empty or $\{2\}$ the table has $t_2<t_3$, so only order 1 is first. There is no fixed $t>0$ at which every $t$ derivative of $g_t(u)$, $u>0$, is nonnegative (otherwise its Taylor series at $t$, of radius $t$, would have a singularity at its positive real radius).
\item Keeping the first factor $e^{-(\log q)^2/4t}$ below $e^{-\gamma_1^2t}$ takes every prime power up to about $e^{2\gamma_1t}$ in this balance, and each atom $\log q$ extends the order one window to about $\log q/2\gamma_1$. These are heuristic scales for this truncation scheme, not theorems of necessary information (a tail can sometimes be bounded without listing its atoms); the full condition on $(0,\infty)$ would need them all.
\end{enumerate}



\begin{proposition}[One sided certificates]\label{wt:prop:cert}
Let $q$ be the smallest prime power omitted from $\mathcal P_S$, so that $W=(\mathcal A-\mathcal P_S)-\mathcal P_{\mathrm{rest}}$, where $\mathcal P_{\mathrm{rest}}$ collects the omitted prime powers, all at least $q$.
\begin{enumerate}
\item[(1)] For $t<(\log q)^2/2$, $-W'(t)\ge-(\mathcal A-\mathcal P_S)'(t)$. Hence if $\mathcal A-\mathcal P_S$ is decreasing on $(0,X]$ with $X<(\log q)^2/2$, so is $W$, with no tail estimate needed.
\item[(2)] For even orders a tail bound is needed. Let $S$ be all prime powers up to $N_0$ and let $f(y)=y^{-1/2}\partial_t^2g_t(\log y)$ be nonnegative and decreasing on $(N_0,\infty)$. Then, with $\psi_{\mathrm C}(y)<1.03883\,y$ \cite{RS62},
\[
W''(t)\ge(\mathcal A-\mathcal P_{\le N_0})''(t)-1.03883\Big(N_0f(N_0)+\int_{\log N_0}^\infty e^{u/2}\partial_t^2g_t(u)\,du\Big).
\]
The hypotheses on $f$ hold for every $0<t\le X$ as soon as $X\le(\log N_0)^2/[2(5+\sqrt{10})]$.
\end{enumerate}
\end{proposition}

\begin{proof}
For (1), $\partial_t\log g_t(u)=u^2/4t^2-1/2t>0$ when $t<u^2/2$, so every term of $\mathcal P_{\mathrm{rest}}$ has positive $t$ derivative, and the differentiated prime sum converges locally uniformly. For (2), Stieltjes integration by parts gives for $f\ge0$ decreasing and integrable with $yf(y)\to0$ that $\sum_{m>N_0}\Lambda(m)f(m)=-\psi_{\mathrm C}(N_0)f(N_0)-\int_{N_0}^\infty\psi_{\mathrm C}f'\,dy\le1.03883\big(N_0f(N_0)+\int_{N_0}^\infty f\,dy\big)$, and $\int_{N_0}^\infty f(y)\,dy=\int_{\log N_0}^\infty e^{u/2}\partial_t^2g_t(u)\,du$. For the hypotheses put $z=(\log y)^2/4t$ and $Q(z)=z^2-3z+\frac34$, so that $f(y)=y^{-1/2}g_t(\log y)t^{-2}Q(z)$; the $u$ derivative of $f$ has the sign of $\frac u{2t}(Q'(z)-Q(z))-\frac12Q(z)$ with $Q'-Q=-z^2+5z-\frac{15}4$, and both $f\ge0$ and $f'\le0$ hold when $z\ge(5+\sqrt{10})/2$.
\end{proof}

\begin{resultpar}\label{wt:res:cert}
We evaluated the lower bound in (2) on the geometric grid $0.01\cdot1.08^k$ of $t$, at 30 digits, checking that $f$ is nonnegative and decreasing on a grid of $u$ of step $0.01$ (computed). The bound stays positive up to the points of Table~\ref{wt:tab:cert}; the first failing grid points, where it is negative, are $0.0799$, $0.1597$ and $0.2534$.\end{resultpar}

\begin{table}[h]\centering\small
\fitwidth{\begin{tabular}{cccc}
\toprule
prime powers kept & order 2 window from the bound & law $\log q/2\gamma_1$ & order 1 window (Table~\ref{wt:tab:pbp})\\
\midrule
$\le 10$ & 0.0740 & 0.0848 ($q=11$) & 0.0856\\
$\le 100$ & 0.1479 & 0.1633 ($q=101$) & 0.1635\\
$\le 1000$ & 0.2346 & 0.2447 ($q=1009$) &\\
\bottomrule
\end{tabular}}
\caption{Windows for $W''>0$ suggested by the bound of Proposition~\ref{wt:prop:cert}(2) from the primes up to $N_0$ with the Rosser and Schoenfeld tail (computed in floating point on a grid; not certified).}\label{wt:tab:cert}
\end{table}

The windows follow the same law and fall short of it by $12.8\%$, $9.4\%$ and $4.1\%$, the cost of the explicit tail. The inequalities are proved; the table is not a certificate, because floating point positivity at grid points proves neither the signs of the computed quantities nor positivity between the points, and it does not cover $0<t<0.01$. A continuous certificate would partition $[\varepsilon,X]$ into intervals $I_i=[a_i,b_i]$, enclose the lower bound $B_N(c_i)\ge L_i$ at a point $c_i\in I_i$ and a bound $|B_N'|\le M_i$ on $I_i$, and check $L_i>M_i\max(c_i-a_i,b_i-c_i)$; below $\varepsilon$, Theorem~\ref{thm:ch12:cmint}(a) covers $(0,0.019]$ for every order. For every order and every $t$ the theorem is Theorem~\ref{thm:ch12:window}, which uses the zeros. On the method of Theorem~\ref{thm:ch12:cmint}: its bound for part (a) vanishes at $t=0.0193$ (Section~\ref{sec:ch12:allorders}), and a reader of that section reported that the sharp envelope of the prime terms in place of the proved one moves the root to about $0.022$ ($0.0218$ with the supremum over $k\le300$; the book does not reproduce it, unchecked), the limit then being set by $E_1$, which vanishes at $0.0261$.
